% chapters/conclusions.tex
\chapter{Conclusions}
\label{ch:conclusions}

\section{Further Work}
\label{sec:6.1}

The future work based upon this project is divided into three categories: use of the existing output as preliminary data, improvements to the current method and addition of complexity to the model, all of which may contribute to the resolution of questions posed by the current results. The first category involves the computation of the Bayes factor for each of the models. This enables a comparison of the global likelihood of two models $M_i , M_j$ which takes into account their different complexity (in this case, the number of parameters varying from the concordance model):

% Equation (6.1)
\begin{equation}
    O_{ij} 
    = \frac{ \prob{M_{i}}[I] }{ \prob{M_{j}}[I] } \frac{ \prob{D}[M_{i},I] }{ \prob{D}[M_{j},I] }
    = \frac{ \prob{M_{i}}[I] }{ \prob{M_{j}}[I] } \frac{ \mathcal{L} (M_{i}) }{ \mathcal{L} (M_{j}) }
    = \frac{ \prob{M_{i}}[I] }{ \prob{M_{j}}[I] } B_{ij}
\label{eq:bayes-factor}
\end{equation}

where the odds factor $O_{ij}$ reduces to the Bayes factor $B_{ij}$ because there is no prior reason to favour one model over the other \cite{Gregory2005}. The global likelihood of the model $M_i$ is an integral over the product of prior $\prob{x}[M_i , I]$ and likelihood $\prob{D}[x, M_i , I]$ (which is obtained as part of the calculation of the posterior):

% Equation (6.2)
\begin{equation}
    \prob{D}[M_{i},I]
    = \int{} \dd{x} \prob{x}[M_{i},I] \prob{D}[x,M_{i},I]
\label{eq:model-likelihood-integral}
\end{equation}

There are several existing methods to calculate tha Bayes factor, all of which rely either on brute-force integration or making an approximation so the the integral can be replaced by the solution of a matrix equation (e.g. the linear least-squares approximation) \cite{Gregory2005}. A new method suggested by Weinberg calculates $B_{ij}$ directly from the posteriors without the need for either method \cite{Weinberg2009}.

The second category (which will require more computing power) increases the complexity of the algorithm used to generate the Markov chains in order to achieve faster convergence to the posterior. This involves either changing the step size using parallel tempering or minimising the chain length by applying a convergence test.

The third category involves extending the complexity of the model: either introducing more parameters and/or allowing more of them to vary at once (once again, the latter will require more computing power). Of these options, the one which is most cosmologically interesting is the addition of more parameters, in particular including a more general form of the dark energy equation of state. The number of cosmological parameters which can be extrapolated from SNe data is small compared to the richness of the \cmb{} \cite{Francis2008} and the majority of additions are the inclusion of more cosmological fluids \cite{Hobson2006}:
\begin{itemize}
    \item a radiation term (dominated by the photons from the \cmb{} which is well-quantified, but also neutrinos which have the same equation of state $w = 1/3$)
    \item the separation of matter into dark and baryonic components (which would allow separate priors from different data)
    \item and a curvature term (not strictly a fluid, but it can be modelled as one) which would allow non-flat geometries.
\end{itemize}
Overall, it is likely that the addition of \cmb{} photons and a more general $w(z)$ would have the greatest impact. The latter two approaches involve repeating the project: this enables the inclusion of the most up to date SNe data (currently \cite{Amanullah2010}, or this may include active observing for Ia SNe and reanalysis of the ”global” SNe data using new light curves). The method from this project (or some improved variant) can also be repeated with data from other sources, notably \cmb{}, \bao{} and constraints from nucleosynthesis. The future work for this project is particularly rich, encompassing numerical, theoretical and observational fields.

\section{Summary}
\label{sec:6.2}

Distance-redshift relations are sensitive to the values of the cosmmological parameters both locally and at high redshift. Distance modulus data obtained from three recent Type Ia SNe surveys was used as a comparative tool to evaluate the consistency of various cosmological models. Five additional SNe sets were artificially generated: three examining how the error in the cosmological parameters scales with reduction of errors from the SNe data and two providing projected results from a newly- proposed space-based survey. The aim of the project was to determine best-fit values and confidence limits for the cosmological parameters across all data sets by application of Bayesian inference, using Markov chain Monte Carlo methods to determine the posteriors from which these values were extracted.

During the course of the project, I demonstrated the following:

\begin{itemize}
\item On large scales, the contents of the Universe are well-approximated by a series of cosmological fluids with various equations of state. If the energy pressure of a fluid is negative, then it falls into a class of fluids known as dark energy.
\item The Friedmann-Lemaitre-Robertson-Walker metric which determines the evolution of the Universe follows directly from the application of the cosmological principle (homogeneity and isotropy) to the field equations of general relativity.
 \item The Friedmann equations are a set of numerically integrable, coupled, first-order ODEs which arise from the FLRW metric and the requirement of local energy conservation. The diagonal form of the metric permits only perfect fluids to be included in the Friedmann equations.
\item This restriction narrows the classes of dark energy which could be examined by this project. In particular it eliminates quintessence fields while still permitting barotropic fluids and a cosmological constant.
\item The luminosity distance-redshift relation is a robust method for estimating values of the cosmological parameters using suitably calibrated Type Ia supernovae as standard candles.
\item Cosmological parameter estimation can be achieved by application of Bayesian inference, which generates a posterior probability density function (PDF) taking into account prior information and observational data.
\item Calculation of the posterior PDF can be achieved numerically with reasonable efficiency using Markov chain Monte Carlo methods, the simplest of which is the Metropolis-Hastings algorithm. The simplicity of this algorithm means that some fine-tuning is required before convergence to the posterior is achieved.
\item The one- and two-parameter posteriors can be converted into credible regions which summarise the probability that a subset of parameter values contains the values corresponding to the observable Universe.
\item Analysis of the one-parameter results demonstrates that the 1d parameter values agree with the concordance model if all data sets are taken into account. Individual data sets show better fits to non-concordance values, however the tension between the Gold set and the Constitution and Union sets results in no single non-concordance values being a good fit to all three SNe sets. 
\item This tension increases when one scales the errors on the observed data sets, demonstrating the difference in light curve calibration between the Gold set and the Constitution/Union sets. When these smaller-error sets are compared to the projected SNAP sets, the uncertainty in the best-fit cosmological parameters is comparable to the uncertainty from the estimated ”pessimistic” result from the proposed SNAP/WFIRST survey. The optimistically estimated result has uncertainties which are superior to current results by up to an order of magnitude.
\item The two-parameter models show degeneracy between parameters, which increases the difficulty in constraining parameter values. Current observation have phase space regions of linear degeneracy between all parameters. Excluding the $(\Omega_m , \gamma)$ and $(w_0 , w_1)$ models, these credible regions curve into areas of non-degeneracy, inside of which a value of one parameter provides no constraints on the others. Examination of the corresponding regions for the artificially generated sets suggests that linear degeneracy will emerge for all of the two-parameter models if the errors in the data are reduced. 
\item Until improved light-curve fitters (\cref{app:type-ia-standard-candles}) are developed, values of the cosmological parameters derived from Type Ia SNe will be highly sensitive to the fitters used. This accounts for the overwhelming majority of the tension between the Gold set and the Constitution/Union sets.
\item These degeneracies force us to obtain data with complementary degeneracies in order to tightly constrain the parameters. This is currently available for a three-parameter space $(\Omega_m , w_0 , w_1 )$ from the seven-year WMAP data from the cosmic microwave background and Sloan Digital Sky Survey data on baryon acoustic oscillations.
\end{itemize}
